\cotefile{157-5}
\scanpage{543}
\mspage{\deleted{4\illegible{1}}454}

soit cofinal dans $J$, donc dans $I$. Comme cet ens \inserted{$J'$} est contenu dans $I_{i',v}$, cela implique que $I_{i',v}$ est cofinal dans $I$, i.e. $v\in N_{i'}$. On a d'autre part, pour $j\in\deleted{\unsure{I_{i,u}|I}}J'$
\[
\underset{v}{\lfloor\!\underline{\alpha_{ji'}}\!\rfloor}\,\alpha_{i'i}=\underset{u}{\lfloor\!\underline{\alpha_{ji}}\!\rfloor}
\]
d'où $u=v\alpha_{i'i}$ avec $v\in N_{i'}$, d'où $u\in N_{\unsure{i'|i}}\alpha_{i'i}$, qed. $\big]$

\deleted{Je dis que $\unsure{R|X}$} \inserted{Pour $i\in I$, soit}
\[
\alpha_i\colon\underbrace{\mathrm{Hom}(\deleted{\unsure{x|a}}\xi\deleted{\unsure{\,\!,|i}},\underset{\substack{\Vert\\\xi}}{x_i})}_{=M}\longrightarrow
\underset{\substack{\Vert\\\underset{i}{\varinjlim}\,\mathrm{Hom}(\xi,x_i)}}{\underline{\mathrm{Hom}(\deleted{\unsure{x|a}}\xi,F)=F(\xi)}}
\]
le morphisme canonique.

On trouve ceci :

\noindent\makebox[0pt][r]{(4.15.8.4)\quad}Si $u\in N_i$, alors la relation d'équivalence $R_i$ dans $M$ définie par $\alpha_i$ ($v\overset{R_i}{\sim}v'\iff\alpha_i(\unsure{v|w})=\alpha_i(v')$) \deleted{\unsure{qui s'égale|équivaut}} \inserted{est égale} à celle définie par l'application $v\longmapsto uv$ de $M$ dans $M$, i.e.
\[
\alpha_i(v)=\alpha_i(v')\quad\Longleftrightarrow\quad uv=uv'.
\]

C'est immédiat : puisque $I_{i,u}$ \unsure{est} cofinal dans $I$, \deleted{\unsure{l'|il}} on a $\alpha_i(v)=\alpha_i(v')$ sss $\exists\, j\in I_{i,u}$ avec $\alpha_{ji}v=\alpha_{j\unsure{i|{}}}v'$, or $\alpha_{ji}=u$.

\deleted{\unsure{Notons|Posons} $R_u$}

Pour $u\in M$, notons $R_u$ la relation d'équivalence dans $M$ définie par $v\longmapsto uv$. L'ens des relations d'équivalence dans
% archive sheet number: 542
